\documentclass{dit-thesis}
\usepackage[utf8]{inputenc}
\usepackage[T1]{fontenc}
\usepackage{geometry}
\usepackage{graphicx}
\usepackage{float}
\usepackage{hyperref}
\usepackage[english,indonesian]{babel}

\title{Optimasi Pemrosesan Template LaTeX Menggunakan Pendekatan Berbasis AI}
\author{Academic Support Agent}
\date{\today}

\begin{document}
\maketitle

\begin{abstract}

\end{abstract}

\chapter{PENDAHULUAN}


\chapter{TINJAUAN PUSTAKA}


% Original content (commented out):
% % Δυνατές επιλογές στην \documentclass{dit-thesis}:
%  - bsc ή msc ή phd: για πτυχιακή ή μεταπτυχιακή εργασία ή διδακτορική %    διατριβή αντίστοιχα
%  - greek ή english: η γλώσσα στην οποία είναι γραμμένη η εργασία
%  - oneside: παραγωγή κειμένου για εκτύπωση μονής όψης

%Για την ανάπτυξη του παρόντος template χρησιμοποιήθηκε το template της πτυχιακής εργασίας της Αγγελικής Νικολακοπούλου με τίτλο "Υποστηρικτικό υλικό για την εκπόνηση πτυχιακών εργασιών"  του Τμήματος Πληροφορικής και Τηλεπικοινωνιών του Πανεπιστημίου Πελοποννήσου.

\documentclass[bsc,greek]{dit-thesis}
% Το πακέτο αυτό παράγει το τυχαίο κείμενο που έχει εισαχθεί στο παράδειγμα της
% εργασίας ώστε να γεμίσουν κάποιες σελίδες. Προφανώς μπορεί να αφαιρεθεί
% σε μία πραγματική εργασία, μαζί με τις εντολές \lipsum παρακάτω.
\usepackage{lipsum}
\usepackage{xcolor}

%Το πακέτο αυτό χρησιμοποιείται για την αλλαγή του σημείου στίξης της λεζάντας από άνω και κάτω τελεία σε τελεία. 
\usepackage[labelsep=period]{caption}

\usepackage{supertabular, booktabs}
\usepackage{multicol}
\usepackage{tabularx}
\usepackage{longtable}
\usepackage{float}
\usepackage{verbatim}
\usepackage{cleveref}


% Βιβλιογραφία.
% Αλλάξτε την παράμετρο style εφόσον επιθυμείτε διαφορετικό βιβλιογραφικό στυλ.
% Δείτε την τεκμηρίωση του πακέτου biblatex σχετικά

%\usepackage[backend=biber, style=apa, natbib=true, hyperref=false]{biblatex}
\usepackage[backend=biber,style=alphabetic]{biblatex}

% Το αρχείο που περιέχει τη βιβλιογραφία της εργασίας.

\addbibresource{references.bib}

\DeclareBibliographyCategory{cited}
\AtEveryCitekey{\addtocategory{cited}{\thefield{entrykey}}}

%Το πακέτο αυτό χρησιμοποιείται  για την διαμόρφωση του περιεχομένου του πίνακα σχημάτων και πινάκων.

\usepackage{tocbasic}

\DeclareCiteCommand{\citeyear}
    {\usebibmacro{prenote}}
    {\bibhyperref{\printfield{year}}\bibhyperref{\printfield{extrayear}}}
    {\multicitedelim}
     {\usebibmacro{postnote}}


\DeclareTOCStyleEntry[
  entrynumberformat=\entrynumberwithprefix{\figurename},
  dynnumwidth,
  numsep=1em
]{tocline}{figure}
\newcommand\entrynumberwithprefix[2]{#1\enspace#2.\hfill}

\DeclareTOCStyleEntry[
  entrynumberformat=\entrynumberwithprefix{\tablename},
  dynnumwidth,
  numsep=1em
]{tocline}{table}

\newcolumntype{L}[1]{>{\raggedright\let\newline\\\arraybackslash\hspace{0pt}}p{#1}}
\newcolumntype{C}[1]{>{\centering\let\newline\\\arraybackslash\hspace{0pt}}p{#1}}
\newcolumntype{R}[1]{>{\raggedleft\let\newline\\\arraybackslash\hspace{0pt}}p{#1}}



% Ο τίτλος της εργασίας. Μπορείτε να χωρίσετε γραμμές σε σημεία που θέλετε
% χρησιμοποιώντας το «\\».
\title{Εδώ εισάγουμε τον τίτλο της πτυχιακής εργασίας - από το αρχείο thesis.tex}

% Οι συγγραφείς της εργασίας.
% Χρειάζεται μία εντολή \author για κάθε συγγραφέα. Η πρώτη παράμετρος
% είναι το ονοματεπώνυμο και η δεύτερη παράμετρος ο αριθμός μητρώου.
\author{Όνομα Φοιτητή}{Αριθμός Μητρώου}
%\author{<\,Όνομα2 Επώνυμο2\,>}{<\,Αριθμός Μητρώου2\,>}

% Οι επιβλέποντες της εργασίας.
% Χρειάζεται μία εντολή \supervisor για κάθε επιβλέποντα. Η πρώτη παράμετρος
% είναι το ονοματεπώνυμο του συγγραφέα και η δεύτερη η ιδιότητά του.
%\supervisor{<\,Όνομα2 Επώνυμο2\,>}{<\,Ιδιότητα2\,>}
\supervisor{Όνομα επιβλέποντα καθηγητή}{Βαθμίδα καθηγητή}

% Μια προαιρετική εντολή για να καθορίσετε την «ετικέτα» των επιβλεπόντων στο
% εξώφυλλο της εργασίας, αν η προκαθορισμένη δεν είναι σωστή.
%\supervisorlabel{Επιβλέπουσα}
%\supervisorlabel{Επιβλέπουσες}

% Ο μήνας και το έτος παρουσίασης της εργασίας.

\date{Μήνα, έτος}
% Φάκελος με τις εικόνες της εργασίας (αν θέλετε να τις έχετε σε χωριστό
% φάκελο και να μπορείτε να αναφερθείτε με αυτές μόνο με το όνομά τους, χωρίς
% να δίνετε και τον φάκελο στον οποίο βρίσκονται.
% Μπορείτε να προσθέσετε περισσότερους φακέλους στην εντολή, όπως φαίνεται
\graphicspath{
  {figures/}
%  {chapter1/figures}
}

%%%%%%%%%%%%%%%%%%%%%%%%%%%%%%%%%%%%%%%%%%%%%%%%%%%%%%%%%%%%%%%%%%%%%%%%%%%%%%%%

% Στο κείμενο της εργασίας που ακολουθεί, μπορείτε να επέμβετε μόνο στα σημεία που σημειώνονται. Οι υπόλοιπες εντολές αφορούν τη μορφοποίηση της εργασίας και θα πρέπει να μην μεταβληθούν.

\begin{document}

\maketitle
\clearemptydoublepage

\frontmatter
\pagestyle{empty}

% Αφιέρωση. Αν δεν θέλετε να έχετε αφιέρωση, βάλτε σε σχόλια ή διαγράψτε τις επόμενες τρεις γραμμές
\begin{dedication}
Εδώ συμπληρώνεται την αφιέρωση της εργασίας σας.
\end{dedication}

% Ευχαριστίες
\include{acknowledgements}

% Στο αρχείο abstracts.tex γράψτε την περίληψη της εργασίας στα ελληνικά και στα αγγλικά.
\include{abstracts}

\pagestyle{fancy}

\tableofcontents
\addcontentsline{toc}{chapter}{\contentsname}
\clearemptydoublepage

%Κατάλογος σχημάτων
\listoffigures
\addcontentsline{toc}{chapter}{\listfigurename}
\clearemptydoublepage

%Κατάλογος πινάκων
\listoftables
\addcontentsline{toc}{chapter}{\listtablename}
\clearemptydoublepage

%Η αρίθμηση των εικόνων και των πινάκων δεν ακολουθεί την αρίθμηση των κεφαλαίων.
\counterwithout{figure}{chapter}
\counterwithout{table}{chapter}


\mainmatter
%κεντρική στοίχιση περιεχομένου κελιών σε πίνακα 
\newcolumntype{M}[1]{>{\centering\arraybackslash}m{#1}}

\renewcommand{\chaptermark}[1]{\markboth{\thechapter. #1}{}}
\renewcommand{\sectionmark}[1]{\markright{\thesection. #1}}

\newcommand{\tg}[1]{\textgreek{#1}}

% Κάθε κεφάλαιο περιλαμβάνεται σε ξεχωριστό αρχείο, για διευκόλυνση της
% συγγραφής. Περιλάβετε εδώ όλα τα κεφάλαια της εργασίας.
% Η παράμετρος της εντολής είναι το όνομα του αντίστοιχου αρχείου .tex
% (παρατηρήστε ότι δεν χρειάζεται η κατάληξη εδώ).
\include{chapter1}
\include{chapter2}
\include{chapter3}
\include{chapter4}
\include{chapter5}
\include{chapter6}
%\appendix

% Περιλάβετε εδώ τυχόν παραρτήματα της εργασίας.
%\include{appendix1}


\backmatter

\fancyhead[LO]{\small\nouppercase{\rightmark}}
\fancyhead[RE]{\small\nouppercase{\leftmark}}


% Επειδή οι περισσότερες βιβλιογραφικές αναφορές θα είναι στα Αγγλικά,
% αλλάζουμε τη γλώσσα
%\begin{english}

\printbibliography[category=cited,heading=bibintoc,title={Βιβλιογραφικές Αναφορές}]

\clearemptydoublepage
%\end{english}


\nocite{*}%
\printbibliography[notcategory=cited, heading=bibintoc,title=Πρόσθετη Βιβλιογραφία]
    
\clearemptydoublepage 

\end{document}

\end{document}